\documentclass[11pt,a4paper]{article}
\usepackage[margin=27mm]{geometry}
\usepackage{amsmath,amssymb,amsthm,fontspec,unicode-math,hyperref,booktabs,longtable}
\setmainfont{STIXGeneral}
\setmathfont{STIX Two Math}
\setmonofont{Menlo}
\usepackage{newunicodechar}
\newunicodechar{∎}{\ensuremath{\blacksquare}}
\hypersetup{hidelinks}
\setlength{\emergencystretch}{2em}
\providecommand{\tightlist}{\setlength{\itemsep}{0pt}\setlength{\parskip}{0pt}}
\newcommand{\gen}[1]{\langle #1\rangle}
\newcommand{\Sym}{\operatorname{Sym}}
\newcommand{\Alt}{\operatorname{Alt}}
\newcommand{\core}{\operatorname{core}}
\newcommand{\F}{\mathbb F}
\newtheorem{theorem}{Theorem}
\newtheorem{lemma}[theorem]{Lemma}
\newtheorem{proposition}[theorem]{Proposition}
\newtheorem{corollary}[theorem]{Corollary}

\title{Inductively minimal geometries and their tree coordinates}
\author{Human Compression project}
\date{}
\begin{document}
\maketitle
\begin{abstract}
The four property conclusions for inductively minimal geometries are obtained from one weighted-forest construction. The proof identifies the precise two-block overgroups, reconstructs the underlying tree from the diagram, and separates ordinary, residual and quasi-primitivity. The foundational classification and uniqueness theorem of Buekenhout, Cara and Dehon remain explicitly retained.
\end{abstract}
\section{Introduction}
Buekenhout and Cara proved that inductively minimal pairs satisfy several strong local and global transitivity properties. We recover their complete four-part theorem by passing through the actual edge cuts of a tree. The common paper [UB] supplies the forest geometry and its stabilizer intervals; all additional diagram and parity arguments are proved here. The available source is the complete authored browser text of the 1998 paper, whose formula and diagram rendering is imperfect. No claim of identity with an unavailable original PDF or of a measured source length is made.
\section{Retained classification and conventions}
For the finite groups considered here, a subgroup is \textbf{quasi-maximal} if the interval from it to the ambient group in the subgroup lattice is a chain, equivalently if it has a unique saturated subgroup chain to that group; QPRI requires this for every element stabilizer. Property \((2T)_1\) means that each panel stabilizer is two-transitive on the rank-one residue. Property \((IP)_2\) means that each rank-two residue is either a complete bipartite graph or has at most one common neighbor for each pair of distinct same-type vertices. This last condition concerns geometric incidence, not the subgroup intersection property of [UB].

The diagram has one vertex for each type; two vertices are joined when their rank-two residue is not complete bipartite. An \textbf{end} of a connected diagram is a vertex whose deletion leaves a connected graph. A \textbf{middle} is a vertex whose deletion leaves at least two components, all of equal cardinality. This spells out the source's nonend requirement; in particular a one-vertex diagram has no middle. A finite firm flag-transitive pair of rank \(r\) and connected diagram is \textbf{minimal} if \(|G|\leq(r+1)!\), and \textbf{inductively minimal} if the pair at every residue with connected diagram is minimal, including the empty flag.

\textbf{B97 (foundational classification).} Buekenhout--Cara--Dehon, \emph{Inductively minimal flag-transitive geometries} (1997), as recalled in {[}BC98{]} Lemma 1 and the following paragraph: an inductively minimal pair of rank \(r\) is thin and residually connected and has group \(S_{r+1}\). Its connected diagram has no induced cycle longer than three, each edge lies in a unique maximal clique, and each vertex lies in one or two maximal cliques. Conversely every such connected diagram determines exactly one inductively minimal pair up to isomorphism. We retain the classification and uniqueness, not the four targeted property conclusions of the 1998 paper.
\section{The unit-weight property profile}
\textbf{Proposition B1.Profile (unweighted tree property profile).} For a tree \(T\) on \(N\geq2\) vertices, the geometry of Theorem UB.F with unit blocks has rank \(N-1\), group \(S_N\), and diagram \(L(T)\), the line graph of \(T\). It is inductively minimal, satisfies RWPRI, \((2T)_1\), and \((IP)_2\), and has the following exact properties:

\begin{itemize}
\tightlist
\item
  PRI holds if no edge bisects \(N\), with the additional primitive case \(N=2\).
\item
  Every element stabilizer is quasi-maximal. Failure of PRI is equivalent to a middle in \(L(T)\).
\item
  RPRI holds exactly when \(T\) is a star, equivalently when \(L(T)\) is complete.
\end{itemize}

\emph{Proof.} The unit-weight specialization of Theorem UB.F gives trivial chamber stabilizer and two elements in each panel; the panel group is \(S_2\), proving thinness and \((2T)_1\). A two-edge forest consists either of a three-vertex path or of two disjoint edges. The corresponding rank-two residue is respectively a six-cycle or a complete bipartite graph \(K_{2,2}\). Consequently the diagram is \(L(T)\) and \((IP)_2\) holds.

Each connected remaining edge set with \(m\) edges is a subtree with \(m+1\) vertices; its residue group is \(S_{m+1}\), while all other vertices are isolated and contribute trivial factors. This verifies inductive minimality. In any nonempty forest residue choose a leaf of a nontrivial component. Its edge cut gives a point stabilizer in that component's symmetric group. This is maximal, and multiplying by the other factors preserves maximality. Thus RWPRI holds.

Lemma UB.Overgroups gives the global PRI and quasi-maximality tests. Removing the vertex \(e\) from \(L(T)\) leaves the line graphs of the two sides of \(T-e\), with respective cardinalities \(a-1\) and \(N-a-1\); an empty side of this diagram corresponds to a leaf cut. A middle therefore exists exactly when an edge has \(a=N-a\geq2\), which is exactly failure of PRI. For \(N=2\) there is no middle and the pair is primitive.

Every nontrivial subtree of a star is again a star, and every edge cut of a star is a leaf cut, so its residues are PRI. A nonstar tree contains a three-edge path. The residue on those three edges is an \(S_4\) geometry with a \(2+2\) middle cut; its stabilizer \(S_2\times S_2\) is not maximal by Lemma UB.Overgroups. Thus RPRI fails. Finally, a tree has pairwise adjacent edges exactly when it is a star, proving the line-graph formulation. ∎
\section{Reconstructing the tree from the diagram}
\textbf{Lemma B1.Reconstruction (diagram recognition).} Every connected graph satisfying the three diagram conditions in B97 is the line graph of a tree, unique up to isomorphism. Hence B97 identifies every inductively minimal pair with one of Proposition B1.Profile.

\emph{Proof.} For a diagram with at least two vertices, make a vertex for each maximal clique. A diagram vertex in two maximal cliques gives an edge between the corresponding clique vertices; a diagram vertex in one maximal clique gives an edge from its clique vertex to a new leaf. Two distinct maximal cliques meet in at most one vertex, since a shared edge would violate the unique-clique condition. Thus the resulting graph has no loops or parallel edges. It is connected, and its edge labels are the original diagram vertices, with exactly their original adjacencies.

A shortest cycle of length at least four in this graph would give an induced cycle of that length on its edge labels in the diagram. A triangle would give three pairwise adjacent edge labels, which extend to a maximal clique; its three pairwise edges would then violate uniqueness of the maximal cliques already represented at the three cycle vertices. Both cases are impossible. The constructed graph is therefore a tree. Conversely the maximal cliques in the line graph of a tree are its nontrivial edge stars, so this construction recovers the tree. A one-vertex diagram corresponds separately to the one-edge tree. Proposition B1.Profile constructs an inductively minimal pair for every such diagram; the uniqueness in B97 identifies it with the given pair. ∎
\section{All four property conclusions}
\textbf{Corollary B1.Source (inductively minimal pairs).} For every finite inductively minimal pair of positive rank:

\begin{enumerate}
\def\labelenumi{\arabic{enumi}.}
\tightlist
\item
  RWPRI, \((2T)_1\) and \((IP)_2\) hold.
\item
  Even rank implies PRI.
\item
  Every type stabilizer is quasi-maximal; failure of PRI is equivalent to a middle in the connected diagram.
\item
  RPRI holds if and only if the diagram is complete.
\end{enumerate}

\emph{Proof.} The retained B97 classification and B1.Reconstruction identify the pair with a unit-weight tree geometry. B1.Profile gives the first, third and fourth conclusions. For the second, an even number of edges gives an odd number of vertices, so no edge cut is balanced; UB.Overgroups gives maximality of every type stabilizer. These are all four clauses of {[}BC98, Theorem 1{]}. The stronger wording in clause 3 avoids an ambiguity in the source: its phrase ``this is the case'' characterizes failure of PRI, not QPRI alone. The one-vertex diagram has no middle and is included. ∎

\section*{References and retained inputs}
{[}UB{]} \emph{Forest coordinates and relative generation}, common theory in the present volume. Its full proofs and definitions are part of this application's mathematical dependencies.

{[}BC98{]} F. Buekenhout and P. Cara, \emph{Some properties of inductively minimal geometries}, Bulletin of the Belgian Mathematical Society--Simon Stevin \textbf{5} (1998), 213--219, DOI \url{10.36045/bbms/1103409005}. The consulted complete authored text is the browser-rendered author upload; its diagram and formula rendering is imperfect. No identity with an unacquired PDF is assumed.

{[}BCD97{]} F. Buekenhout, P. Cara and M. Dehon, \emph{Inductively minimal flag-transitive geometries}, in N. L. Johnson (ed.), \emph{Mostly Finite Geometries}, Marcel Dekker, 1997, 185--190. B97 is retained in the exact form recalled in {[}BC98, Lemma 1 and the following paragraph{]}; the original chapter's proof is not reproduced or claimed independently read here.

\end{document}
